\begin{definition}{Rationals ($\mathbb{Q}$)}{rationals}
    A rational number is an expression of the form $a \quot b$, where $a$ and $b$ are integers and $b$ is non-zero.
    Two rational numbers are considered equal ($a \quot b = c \quot d$) if and only if $ad = bc.$
\end{definition}

% Problem 1
\begin{problem}{Equality on the rationals follow the three axioms}
	Verify that the definition of equality on the rationals is reflexive, symmetric, and transitive.
\end{problem}
\begin{proof}
    First, we show that the equality is reflexive. 
    Consider the rational $a \quot b$.
    Notice that $ab = ba$. 
    Hence, $a \quot b = a \quot b$.

    Second, we prove that the equality is symmetric. 
    Suppose $a \quot b = a' \quot b',$ so $ab' = a'b.$
    We can reverse the right and left hand sides to get $a'b = ab'$.
    Then, we use the commutativity of multiplication for the integers to get $a'b = b'a$, so $a' \quot b' = a \quot b.$
\end{proof}

% Problem 2
\begin{problem}{Addition, multiplication, and negation are well-defined}
    Let $a \quot b, c \quot d \in \mathbb{Q}$. 
    We define addition as $a \quot b + c \quot d = (ad + bc) \quot (bd)$;
    multiplication as $(a \quot b) \times (c \quot d) = (ac) \quot (bd)$;
    and negation as $-(a \quot b)=(-a) \quot b$.
\end{problem}
\begin{proposition}
    Addition on the rationals is well-defined.
\end{proposition}
\begin{proof}
    This has been done in the book.
\end{proof}

\begin{proposition}
    Multiplication on the rationals is well-defined.
\end{proposition}
\begin{proof}
    Suppose $a \quot b = a' \quot b',$ so $a b' = ba'.$
    We show that $(a \quot b) \times (c \quot d) = (a' \quot b') \times (c \quot d)$.
    By definition, the LHS is $(ac) \quot (bd)$, and the RHS is $(a'c) \quot (b'd).$
    Notice:
    \begin{align*}
        (a \quot b) \times (c \quot d) = (a' \quot b') \times (c \quot d) &\iff \big((ac) \quot (bd) \big)= \big((a'c) \quot (b'd) \big) &&\text{by definition of multiplication}\\
        &\iff (ac)(b'd) = (bd)(a'c) && \text{by definition of equality}\\
        &\iff acb' = ba'c && \text{by cancellation law} \\
    \end{align*}
    We now consider two cases: $c = 0$ or $c \neq 0.$ 
    \begin{enumerate}
        \item Say $c = 0.$ Then $acb' = ba'c$ implies $0 = 0.$
        \item Say instead that $c \neq 0.$ Then, by the cancellation law, $ab'=ba'$. This is true by our supposition.
    \end{enumerate}
    In either case, we see that the equality of $\big((ac) \quot (bd) \big)= \big((a'c) \quot (b'd) \big)$ holds
    because the reversible rearrangements lead to $ab' = ba'$, which is true by supposition.
    The proof holds similarly for the case of $c \quot d = c' \quot d'.$

\end{proof}

\begin{proposition}
    Negation on the rationals is well-defined.
\end{proposition}
\begin{proof}
    Suppose $a \quot b = a' \quot b',$ so $ab' = ba'.$
    We show that $-(a \quot b) = -(a' \quot b').$
    By definition of negation, the LHS is $(-a) \quot b$ while the RHS is $(-a') \quot b'.$ Notice:
    \begin{align*}
        -\big(a \quot b \big)= -\big(a' \quot b' \big)  &\iff (-a) \quot b = (-a') \quot b' &&\text{by definition of negation}\\
        &\iff (-a)b' = b(-a') && \text{by definition of multiplication}\\
        &\iff ab' = ba' && \text{by cancellation law} \\
    \end{align*}   
    In either case, we see that the equality of $-\big(a \quot b \big)= -\big(a' \quot b' \big)$ holds
    because the reversible rearrangements lead to $ab' = ba'$, which is true by supposition.
\end{proof}

\break
% Problem 3
\begin{problem}{Laws of Algebra for the Rationals}
    Let $x = x_1 \quot x_2, y = y_1 \quot y_2, z = z_1 \quot z_2$ be rationals where $x_1, x_2, y_1, y_2, z_1, z_2 \in \mathbb{Z}$. 
    Then, the following laws of algebra hold.
\end{problem}


% have multicolumns here
\begin{multicols}{2}
\begin{proposition}
	$x + y = y + x$
\end{proposition}
\begin{proof}
    Notice:
    \begin{align*}
        x + y &= x_1 \quot x_2 + y_1 \quot y_2 \\
        &= (x_1 y_2 + x_2 y_1) \quot (x_2 y_2) \\
        &= (y_2 x_1 + y_1 x_2) \quot (y_2 x_2) \\
        &= (y_1 x_2 + y_2 x_1) \quot (y_2 x_2) \\
        &= y + x
    \end{align*}
\end{proof}

\begin{proposition}
	$(x+y)+z =x+(y+z)$
\end{proposition}
\begin{proof}
    This has been done in the book.
\end{proof}

\begin{proposition}
	$x+0 = 0+x = x$
\end{proposition}
\begin{proof}
    Notice:
    \begin{align*}
        x + 0 &= (x_1 \quot x_2) + (0 \quot 1) \\
        &= (x_1 1 + x_2 0) \quot (x_2 1) \\
        &= (1 x_1 + 0x_2) \quot (1 x_2) \\
        &= x_1 \quot x_2 \\
        &= x
    \end{align*}
    Because addition is commutative, we have that $0 + x = x.$
\end{proof}

\begin{proposition}
	$x + (-x) = (-x) + x = 0$
\end{proposition}
\begin{proof}
    Notice:
    \begin{align*}
        x + (-x) &= (x_1 \quot x_2) + (-(x_1 \quot x_2)) \\
        &= (x_1 \quot x_2) + ((-x_1) \quot x_2) \\
        &= (x_1 x_2 + x_2 (-x_1)) \quot (x_2 x_2) \\
        &= (x_1 x_2 - x_1 x_2) \quot (x_2 x_2) \\
        &= 0 \quot (x_2 x_2) \\
        &= 0
    \end{align*}
    Because addition is commutative, we have that $(-x) + x = 0.$
\end{proof}

\begin{proposition}
	$xy = yx$
\end{proposition}
\begin{proof}
    Notice:
    \begin{align*}
        xy &= (x_1 \quot x_2) \times (y_1 \quot y_2) \\
        &= (x_1 y_1) \quot (x_2 y_2) \\
        &= (y_1 x_1) \quot (y_2 x_2) \\
        &= yx 
    \end{align*}
\end{proof}

\begin{proposition}
	$(xy)z = x(yz)$.
\end{proposition}
\begin{proof}
    Notice:
	\begin{align*}
        (xy)z &= \big((x_1 \quot x_2) \times (y_1 \quot y_2) \big) \times (z_1 \quot z_2) \\
        &= \big((x_1 y_1) \quot (x_2 y_2) \big) \times (z_1 \quot z_2) \\
        &=  \big((x_1 y_1) z_1 \big)\quot \big((x_2 y_2) z_2 \big)\\
        &= \big(x_1 (y_1 z_1) \quot x_2 (y_2 z_2)\big) \\
        &= (x_1 \quot x_2) \times \big((y_1 z_1) \quot (y_2 z_2) \big) \\
        &= (x_1 \quot x_2) \times \big( (y_1 \quot y_2) \times (z_1 \quot z_2) \big) \\
        &= x(yz)
    \end{align*}
\end{proof}

\begin{proposition}
	$x1 = 1x = x$
\end{proposition}
\begin{proof}
    Notice:
    \begin{align*}
        x1 &= (x_1 \quot x_2) \times (1 \quot 1) \\
        &= (x_1 1) \quot (x_2 1) \\
        &= (1x_1) \quot (1 x_2) \\
        &= (x_1) \quot x_2 \\
        &= x
    \end{align*}
    Because multiplication is commutative, $1x = x$.
\end{proof}
\begin{proposition}
	$x(y+z) = xy +xz$
\end{proposition}
\begin{proof}
    There are diminishing returns if I do this.
\end{proof}

\begin{proposition}
	$(y + z) x = yx + zx$
\end{proposition}
\begin{proof}
	Notice:
	\begin{align*}
		(y + z) x &= x (y + z) \\
		&= xy + xz \\
		&= yx + zx.
	\end{align*}
\end{proof}

\begin{proposition}
    If $x$ is non-zero, we have $xx^{-1} = x^{-1}x = 1$
\end{proposition}
\begin{proof}
    Suppose $x$ is non-zero, so $x = a \quot b$ where $a$ and $b$ are non-zero integers. Notice:
    \begin{align*}
        xx^{-1} &= (x_1 \quot x_2) \times (x_1 \times x_2)^{-1} \\
         &= (x_1 \quot x_2) \times (x_2 \times x_1) \\
         &= (x_1 x_2) \quot (x_2 x_1) \\
         &= (x_1 x_2) \quot (x_1 x_2) \\
         &= 1
    \end{align*}
\end{proof}

\end{multicols}

\break
% Problem 4
\begin{problem}{Trichotomy of $\mathbb{Q}$}
    Let $x \in \mathbb{Q}.$ We show the trichotomy of $\mathbb{Q}$, where exactly one of the following statements hold:
    $x = 0,$ $x$ is a positive rational number, or $x$ is a negative rational number.
\end{problem}
\begin{proof}
    Let $x = a \quot b \in Q$ where $a$ is is an integer and $b$ is a non-zero integer.
    We consider six cases, which are based on the possible combinations of $a$ and $b$ being a positive integer, negative integer, or zero.
    \begin{enumerate}
        \item Say $a < 0$ and $b < 0$, so there exist positive natural numbers $a', b'$ such that $a = -a'$ and $b = -b'$.
        By substitution, $a \quot b = (-a') \quot (-b')$.
        Notice also that $(-a') \quot (-b') = a' \quot b'$ because $(-a')b'=(-b')a'$. 
        Hence, $a \quot b = a' \quot b'$ where the RHS is positive because $a', b' > 0$.  
        Therefore, $a \quot b$ is positive.

        \item Say $a < 0$ while $b > 0$, so there exists a positive natural number $a'$ such that $a = -a'.$
        Hence, $a \quot b = (-a') \quot b = -(a' \quot b).$ 
        Because $a \quot b$ is the negation of $a' \quot b$, $a \quot b$ is negative.
        
        \item Say $a > 0$ and $b < 0$, so there exists a positive natural number $b'$ such that $b = -b'.$ 
        By substitution, $a \quot b = a \quot (-b')$.
        Notice that because $ab' = (-b')(-a)$. we have that $a \quot (-b') = (-a) \quot b'$ by definition of equality.
        Observe that $(-a) \quot b'$ = $-(a \quot b')$.
        Therefore, we have a long list of equalities as summarized: 
        $$a \quot b = a \quot (-b') = (-a) \quot b' = -(a\quot b')$$ 
        This is all to just say that $a \quot b$ is the negation of the positive rational $a \quot b'.$
        Hence, $a \quot b$ is negative.

        \item Say $a > 0$ and $b > 0.$ 
        By the definition of a positive rational, $a \quot b$ is also positive.

        \item Say $a = 0$. 
        Notice that for any $b \in \mathbb{N} \setminus \{0\},$ we have that $a \quot b = 0 \quot b = 0$
        We show this through the equivalence eqaution.
        Notice: 
        \begin{align*}
            0 \quot b = 0 \quot 1 &\iff 0 \times 1 = b \times 0 && \text{by definition of equality} \\
            &\iff 0 = 0 &&\text{by multiplication}
        \end{align*}
    \end{enumerate}

    We have shown that for every rational, they are at least either $x =0$, $x > 0$ or $x < 0.$ 
    However, we now need to show that this is mutually exclusive, so they are also AT MOST $x =0$, $x > 0$ or $x < 0.$ 
    Suppose $x = 0$, so $x = 0 \quot b.$ 
    Hence, $a = 0$, so $a$ is not a positive number. Therefore, $x$ cannot be positive.
    Similarly, suppose for the sake of contradiction that $x = 0$ and $x < 0$. Hence, $x = 0 \quot 1$.
    Also, $x = (-c) \quot d$ for some positive numbers $c, d$. Then, $0 \quot 1 = (-c) \times d$, so $0d = (-c)1 = -c$.
    This implies that $c = 0$, but $c$ is positive. A contradiction!

    Suppose for the sake of contradiction that $x < 0$ and $x > 0$, so 
    $x = a \quot b$ and $x = (-c) \quot d$ for positive numbers $a, b, c, d.$
    Then, $a \quot b = (-c) \quot d$ implies $ad = b(-c) = -(bc)$.
    Since $a, b, c, d$ are positive, $ad$ and $bc$ are positive, and $-(bc)$ is negative. 
    However, since $ad = -(bc)$, we now have that $ad$ is also negative.
    By the trichotomy of integers, $ad$ can only be at most positive, negative, or zero. A contradiction!
    
\end{proof}

These three lemmas are here for my sanity because I really want to show that if you multiply by a positive number, the sign does not change.
However, if you multiply by a negative number, the sign does flip.
I just add these lemmas for completeness and especially because I will be using them for the next problem.

\begin{lemma}\label{lem: preservation_mul_positive}
    Let $x = x_1 \quot x_2$ and $y = y_1 \quot y_2$ be rational numbers, where $x_1, x_2, y_1, y_2$ are integers. 
    If $y$ is positive, then $xy$ has the same sign as $x$.
\end{lemma}
\begin{proof}
    Suppose $y$ is positive, so $y_1$ and $y_2$ are positive integers.
    We break the proof down into three cases of $x$ being positive, zero, or negative.
    \begin{enumerate}
        \item Say $x$ is positive, so $x_1$ and $x_2$ are also positive. 
        Then, $xy = (x_1 y_1) \quot (x_2 y_2)$ is also positive because $x_1 y_1$ and $x_2 y_2$ are positive.
        \item Say $x$ is zero. Then, $xy = 0y = 0.$
        \item Say $x$ is negative, so there exists a positive rational number $x'$ such that $x = -x'.$ 
        Then, $xy = (-x')y= -(x'y).$ Since $x'$ and $y$ are positive, it follows from Case (1) that $x'y$ is positive. 
        Hence, $xy$ is the negation of the positive rational $x'y$.
        Therefore, $xy$ is negative.
    \end{enumerate}
    In the three cases, we have shown that $xy$ has the same sign as $x$.
\end{proof}

\break
\begin{lemma}\label{lem:negation_same_multiplication_negative_one}
    Multiplying by $-1$ is the same as negation for the rationals too (i.e., $\forall x \in \mathbb{Q}, \, (-1) \times x = -x$.)
\end{lemma}

\begin{proof}
    Let $x = x_1 \quot x_2 \in \mathbb{Q}$ where $x_1, x_2 \in \mathbb{Z}.$ Notice:
    \begin{align*}
        (-1) \times x &= (-1 \quot 1) \times (x_1 \quot x_2) && \text{by substitution of $-1$ and $x$}\\
        &= (-1 x_1) \quot (1 x_2) && \text{by definition of multiplication on the rationals}\\
        &= (-x_1) \quot (x_2) && \text{by definition of multiplication on the integers} \\
        &= -(x_1 \quot x_2) && \text{by definition of negation}\\
        &= -x &&\text{by substitution}
    \end{align*}
\end{proof}

\begin{lemma}\label{lem: minus_minus_positive}
    The negation of a negative rational number $-(-x)$ is positive.
\end{lemma}
\begin{proof}
    Let $-x$ be a negative rational number.
    Observe that     $-(-x) = (-1) \times \big((-1) \times x\big) = (-1 \times -1) \times x = 1 \times x = x.$
\end{proof}


% Problem 5
\begin{problem}{Properties of the order on the rationals}
    For the properties of the order on the rationals, let $x, y, z \in \mathbb{Q}.$ 
    Then, the following properties hold.
\end{problem}

\begin{multicols}{2}
\begin{proposition}
    (Order Trichotomy) Exactly one of the three statements hold:
    (i) $x = y$; (ii) $x < y$; (iii) $x > y.$
\end{proposition}

\begin{proof}
    Consider the rational $x - y$. By the law of trichotomy for rationals, $x-y$ can at most be negative, zero, or positive.
    If $x - y$ is negative, then $x < y.$ If $x -y = 0$, then $x = y.$ Finally, if $x -y$ is positive, then $x >y.$
    This is mutually exclusive and collectively exhaustive by the law of trichotomy.
    Hence, exactly one of the statements hold for any rational numbers $x$ and $y$.
\end{proof}

\begin{proposition}
    (Order is anti-symmetric) We have that $x < y$ if and only $y >x.$
\end{proposition}
\begin{proof}
    Suppose $x <y$, so $x - y$ is negative.
    Notice that the negation $-(x - y) = y - x$ is positive by Lemma~\ref{lem: minus_minus_positive}.
    Hence, $y>x$.

    Conversely, suppose that $y > x$, so $y - x$ is negative.
    Notice that the negation $-(y-x)=x - y$ is positive by Lemma~\ref{lem: minus_minus_positive}.
    Hence, $x > y.$
\end{proof}

\begin{proposition}
    (Order is transitive) If $x < y$ and $y < z$, then $x < z.$
\end{proposition}
\begin{proof}
    Suppose $x < y$ and $y < z$, so $x - y$ and $y - z$ is negative. 
    Therefore, there exist positive rational numbers $u$ and $v$ such that $x - y = -u$ and $y - z=-v.$
    If we add $-u$ to $-v$, we get $(-u) + (-v) = (-1 \times u) + (-1 \times v)$ by Lemma~\ref{lem:negation_same_multiplication_negative_one}.
    We then use the distributive property to have that $(-1 \times u) + (-1 \times v) = -1(u + v).$ 
    Overall, we have that $(-u) + (-v) = -(u + v)$.
    Since $u + v$ is a positive number, we have that $(-u) + (-v)$ is a negative number since it is the negation of $u + v$.
    
    Notice also that $(-u) + (-v) = (x - y + y - z)= x - z.$
    Overall, we now know that $x - z$ is a negative number. Therefore, $x < z.$
\end{proof}

\begin{proposition}
    (Addition preserves order) If $x < y$, then $x + z < y + z.$
\end{proposition}
\begin{proof}
    Suppose $x < y$, so $x - y$ is a negative rational number.
    Notice that $(x + z) - (y + z) = x + z - y - z= x -y$ is also a negative rational number.
    Therefore, $x + z < y + z.$
\end{proof}

\begin{proposition}
    (Positive multiplication preserves order) If $x < y$ and $z$ is positive, then $xz < yz.$
\end{proposition}
\begin{proof}
    Suppose $x <y$ and $z$ is postive. 
    Then, so $x - y$ is negative. 
    By Lemma~\ref{lem: preservation_mul_positive}, $(x-y)z = xz - yz$ is still negative.
    Hence, $xz < yz.$
\end{proof}
\end{multicols}

\begin{problem}{Multiplying by a negative number reverses order.}
    Show that if $x, y, z$ are rational numbers such that $x < y$ and $z$ is negative, then $xz > yz.$
\end{problem}
\begin{proof}
    Let $x, y, z$ be rational numbers. 
    Suppose $x < y$ and $z$ is negative.
    By the first supposition, $x -  y$ is negative, so there exists a positive rational $u$ such that $x - y = -u.$
    By the second supposition, there exists a positive rational number $v$ such that $z = -v.$
    Notice that $(-u)(-v) = (-1 \times u) (-1 \times v)$ by Lemma~\ref{lem:negation_same_multiplication_negative_one}.
    Then, we can use the commutativity and associativity of multiplication to get $(-u)(-v) = (-1 \times -1) (uv)=1(uv)=uv.$
    This is a very long process to just say $(-u)(-v) = uv.$
    By Lemma~\ref{lem: preservation_mul_positive}, since $u$ and $v$ are both positive, $uv$ is positive. Hence, $(-u)(-v)$, which is equal to $uv$, is also positive.

    Now, observe also that $(-u)(-v) = (x - y)z = xz - yz$.
    Since we have already previously determined that $(-u)(-v)$ is positive, it follows that $xz - yz$ is positive.
    Therefore, $xz > yz.$
\end{proof}